\section{Einrichtung des Testsystems}\label{appendix1:Einrichtung_des_Testsystems}
Zur Einrichtung des Testsystems wird zuallererst ein Linux-Container erstellt, welcher ein einfaches und schnelles Deployen auf unterschiedlichen Linux-Distributionen ermöglicht. Zum Hosten des Linux-Containers wird in dieser Installationsanleitung Ubuntu 20.04 Focal Fossa verwendet, welches eine LTS-Version von Ubuntu darstellt. Nach der Einrichten des Linux-Containers können die benötigten Pakete in dieser neu erstellten Umgebung installiert und daraufhin entsprechend konfiguriert werden. Der vorkonfigurierte Linux-Container ist auch auf der CD/ISO im Verzeichnis \quoteMark{Linux-Container/} vorzufinden. Dabei ist jedoch zu beachten, dass auf dem Host-System dennoch ein paar wenige Konfigurationsschritte für das Host-System durchgeführt werden müssen. Die Schritte für den Linux-Container können hingegen mit dem vorkonfigurierten Linux-Container übersprungen werden.

\subsection{Erstellen des Linux-Containers}
Für die Installation der Linux-Container-Tools werden Administratorrechte auf dem Zielsystem benötigt. Die benötigten Installationsschritte zum Einrichten des Linux-Containers sind wie folgt \cite{lxd_blog_nodate}.

\begin{enumerate}[label*=\arabic*.]
    \item LXD 4.0 installieren:\
        \begin{lstlisting}[style=bash]
sudo snap install lxd
        \end{lstlisting}
    bzw.
        \begin{lstlisting}[style=bash]
sudo apt install lxd
        \end{lstlisting}

    \item LXD mit den Standardeinstellungen initialisieren:\
        \begin{lstlisting}[style=bash]
lxd init
        \end{lstlisting}

    \item Den derzeitigen User zur LXD-Gruppe hinzufügen:\
        \begin{lstlisting}[style=bash]
sudo adduser $(whoami) lxd
        \end{lstlisting}
        Logout \& Login durchführen oder mittels des folgenden Befehls den User im Terminal neu anmelden.
        \begin{lstlisting}[style=bash]
su - $(whoami)
        \end{lstlisting}

    \item Erstellen des Linux-Containers:\\
        Der Linux-Container wird mit dem folgenden Befehl im Default-Linux-Container-Pool erstellt. Der Platzhalter \quoteMark{<container name>} muss durch einen entsprechenden Namen für den Linux-Container ersetzt werden.
        \begin{lstlisting}[style=bash]
CONTAINER_NAME="<container name>"
lxc init ubuntu:18.04 $CONTAINER_NAME
        \end{lstlisting}
        Mit der Option \quoteMark{-s} kann ein anderer Linux-Container-Pool ausgewählt werden.
        \begin{lstlisting}[style=bash]
lxc init ubuntu:18.04 $CONTAINER_NAME -s <container pool name>
        \end{lstlisting}

    \item Die Eigenschaften des Image des neu erstellten Linux-Containers können mit dem folgenden Befehl eingesehen werden:\
        \begin{lstlisting}[style=bash]
lxc storage list
        \end{lstlisting}

    \item Dem Linux-Container erlauben den X11-Traffic weiterzuleiten:
        \begin{enumerate}[label*=\arabic*.]
            \item Die Datei \quoteMark{gui-profile.yaml} mit dem folgenden Inhalt erstellen:
                \begin{lstlisting}[language=yaml]
config:
  environment.DISPLAY: :0
  raw.idmap: both 1000 1000
  user.user-data: |
    #cloud-config
    package_upgrade: true
    packages:
      - x11-apps
      - mesa-utils
      - pulseaudio
    runcmd:
      - 'sed -i "s/; enable-shm = yes/enable-shm = no/g" /etc/pulse/client.conf'
      - 'echo export PULSE_SERVER=unix:/tmp/.pulse-native | tee --append /home/ubuntu/.profile'
description: GUI LXD profile
devices:
  PASocket:
    path: /tmp/.pulse-native
    source: /run/user/1000/pulse/native
    type: disk
  X0:
    path: /tmp/.X11-unix/X0
    source: /tmp/.X11-unix/X0
    type: disk
  mygpu:
    type: gpu
name: gui
used_by:
                \end{lstlisting}

                \item Bei der Erstellung der Datei \quoteMark{gui-profile.yaml} ist zu beachten, dass die IDs im Eintrag \quoteMark{raw.idmap: both 1000 1000} folgende Bedeutung haben:
                    \begin{lstlisting}[language=yaml]
raw.idmap: both <host user id> <user id in container>
                    \end{lstlisting}

                    \begin{tabular}{|>{\raggedright\arraybackslash}p{0.18\textwidth}|>{\raggedright\arraybackslash}p{0.65\textwidth}|}
                        \hline
                        <host user id>: & Die User-ID, welche zum Starten des Linux-Containers verwendet wird.\\
                        \hline
                        <user id in container>: & Die User-ID des Users innerhalb des Linux-Containers. Normalerweise ist dies die User-ID 1000, da der erste User, welcher im Linux-Container eingerichtet wird, standardmäßig die User-ID 1000 erhält.\\
                        \hline
                    \end{tabular}\\

                \item Das Profil kann daraufhin über die Datei \quoteMark{gui-profile.yaml} geladen werden:
                    \begin{lstlisting}[style=bash]
lxc profile create gui
cat gui-profile.yaml | lxc profile edit gui
                    \end{lstlisting}

                \item Abschließend muss das Profil dem neu erstellten Linux-Container zugewiesen werden:
                    \begin{lstlisting}[style=bash]
lxc profile add $CONTAINER_NAME gui
                    \end{lstlisting}
        \end{enumerate}

% Debugging "Cloud-config" with:
% /var/log/cloud-init-output.log
% lxc console ros-tir --show-log

    \item Den Linux-Container starten und eine Shell im Linux-Container öffnen:
        \begin{lstlisting}[style=bash]
lxc start $CONTAINER_NAME
lxc exec $CONTAINER_NAME -- sudo --user ubuntu --login
        \end{lstlisting}

    \item Um Rendering-Probleme zu vermeiden sollten die neuesten Treiber im Linux-Container installiert werden:
        \begin{lstlisting}[style=bash]
sudo add-apt-repository ppa:oibaf/graphics-drivers -yn
sudo apt upgrade -y
        \end{lstlisting}

    \item Um die durchgeleitete Grafikkarte zu testen muss der Linux-Container neu gestartet werden! Anschließend müssen die folgenden Befehle eingegeben werden:
        \begin{lstlisting}[style=bash]
lspci | grep VGA
        \end{lstlisting}
        bzw.
        \begin{lstlisting}[style=bash]
sudo apt-get install glmark2
glmark2
        \end{lstlisting}

        \begin{redbox}{Wichtig:}
            Anstatt \quoteMark{llvmpipe} sollte die jeweilige verbaute Grafikkarte angezeigt werden.
        \end{redbox}

    \item Für die Soundausgabe sollte PulseAudio getestet werden. Es sollte beim Ausführen des folgenden Befehls ein Rauschen zu hören sein:
        \begin{lstlisting}[style=bash]
pacat < /dev/urandom
        \end{lstlisting}

    \item Von Zeit zu Zeit sollte auf dem Host-System ein Backup vom Linux-Container erstellt werden:
        \begin{lstlisting}[style=bash]
lxc export $CONTAINER_NAME ${CONTAINER_NAME}.tar.gz
        \end{lstlisting}

        bzw. zum Erstellen eines Snapshots

        \begin{lstlisting}[style=bash]
lxc snapshot $CONTAINER_NAME snap01
lxc publish ${CONTAINER_NAME}/snap01 --alias <new-alias>
lxc image export <new-alias> ${CONTAINER_NAME}.tar.gz
        \end{lstlisting}

    \item Das Backup kann auf dem Host-System mittels des folgenden Befehls zurückgespielt werden:
        \begin{redbox}{Wichtig:}
            Beim Import des Linux-Containers muss sichergestellt werden, dass zuvor alle mit dem Linux-Container verknüpften Linux-Container-Profile vorhanden sind, da es ansonsten zu Import-Fehlern kommt.
        \end{redbox}

        \begin{lstlisting}[style=bash]
lxc import <new-alias> ${CONTAINER_NAME}.tar.gz
        \end{lstlisting}

        bzw. der Snapshot

        \begin{lstlisting}[style=bash]
lxc image import ${CONTAINER_NAME}.tar.gz --alias <new-alias>
        \end{lstlisting}
\end{enumerate}


\subsection{Installation \& Konfiguration der Pakete}
Mit der zuvor installierten \quoteMark{Ubuntu 20.04 Focal Fossa}-Umgebung können nun die für diese Arbeit benötigten Pakete innerhalb dieser Umgebung installiert und konfiguriert werden.

\begin{enumerate}[label*=\arabic*.]
    \item Installation von ROS 1:\\
    Zur Nutzung von ROS 1 muss ROS Melodic installiert werden. Die Installationsanleitung kann der offiziellen ROS-Webseite (\href{http://wiki.ros.org/melodic/Installation/Ubuntu}{http://wiki.ros.org/melodic/Installation/Ubuntu}) entnommen werden. Bei der Installation des ROS-Pakets muss das \quoteMark{ros-melodic-desktop-full}-Paket mit dem folgenden Befehl installiert werden.

        \begin{lstlisting}[style=bash]
sudo apt install ros-melodic-desktop-full
        \end{lstlisting}

    \item Installation des WidowX 200:
        \begin{enumerate}[label*=\arabic*.]
            \item Hierzu müssen zuvor die benötigten Abhängigkeiten installiert werden:

                \begin{lstlisting}[style=bash]
sudo apt install ros-melodic-robot-state-publisher ros-\
melodic-rqt-plot ros-melodic-joy
                \end{lstlisting}

            \item Anschließend müssen die auf der Webseite \href{https://github.com/Interbotix/interbotix_ros_arms#quickstart}{https://github.com/Interbotix/interbotix\\\_ros\_arms\#quickstart} beschriebenen Schritte 1 bis 6 durchgeführt werden:

            \begin{redbox}{Wichtig:}
                Anstelle des \quoteMark{https://github.com/Interbotix/interbotix\_ros\_arms.git}-Repositories muss, dass für diese Arbeit modifizierte Repository, welches auf der CD/ISO im Verzeichnis \quoteMark{Repositories/interbotix\_ros\_arms} beigelegt ist, verwendet werden. Zum Klonen des modifizierten Repositories muss hierzu der folgende Befehl verwendet werden:

                \begin{lstlisting}[style=bash]
git clone <path to repository on cd or iso/Repositories/interbotix\
_ros_arms>
                \end{lstlisting}
            \end{redbox}

            \item Danach müssen die folgenden Einstellungen auf dem Host-System durchgeführt werden, damit der unprivilegierte Linux-Container Zugriff auf den Dynamixel U2D2 erhält:

                \begin{lstlisting}[style=bash]
wget https://raw.githubusercontent.com/Interbotix/interbotix_ros_arms/\
master/interbotix_sdk/10-interbotix-udev.rules
sudo cp 10-interbotix-udev.rules /etc/udev/rules.d/
sudo udevadm control --reload-rules && sudo udevadm trigger
                \end{lstlisting}

            \item Außerhalb des Linux-Containers müssen folgende Einstellungen erfolgen, damit der unprivilegierte Linux-Container Zugriff auf den Dynamixel U2D2 erhält:
                \begin{lstlisting}[style=bash]
lxc config device add $CONTAINER_NAME dynamixel_u2d2_ttydxl unix-char path\
=/dev/ttyDXL required=false
                \end{lstlisting}

            \item Innerhalb des Linux-Containers muss zudem der folgende Befehl ausgeführt werden, um als normaler User auf den Dynamixel U2D2 zugreifen zu können:
                \begin{lstlisting}[style=bash]
echo "# Interbotix WidowX 200
sudo chown ubuntu /dev/ttyDXL" >> ~/.bashrc
source ~/.bashrc
                \end{lstlisting}
        \end{enumerate}

    \item Testen des WidowX 200:
        \begin{redbox}{Wichtig:}
            \begin{compactitem}
                \item Genügend Abstand einhalten
                \item Den Roboter festschrauben
                \item Beim Beenden der Softwaretools den Roboter festhalten, da dieser ansonsten in sich zusammenfällt.
            \end{compactitem}
        \end{redbox}

        Zum Testen der grundlegenden Funktionalität des WidowX 200 können ein paar spezifische ROS-Befehle an diesen gesendet werden.

        \begin{lstlisting}[style=bash]
roslaunch interbotix_descriptions description.launch robot_name:=wx200 jnt_pub\
_gui:=true

roslaunch interbotix_sdk arm_run.launch robot_name:=wx200

rosservice call /wx200/torque_joints_on

roslaunch interbotix_gazebo gazebo.launch robot_name:=wx200
rosservice call /gazebo/unpause_physics
        \end{lstlisting}

%     \item Wichtige Informationen für ROS bekanntmachen:
%         \begin{lstlisting}[style=bash]
% echo "# ROS
% source /opt/ros/melodic/setup.bash
% source ~/interbotix_ws/devel/setup.bash" >> ~/.bashrc
% source ~/.bashrc
%         \end{lstlisting}

    \item Azure Kinect installieren:
        \begin{enumerate}[label*=\arabic*.]
            \item Außerhalb des Linux-Containers müssen die folgenden Einstellungen erfolgen um die Azure Kinect für den derzeitgen User verfügbar zu machen \cite{microsoftazure-kinect-sensor-sdk_installation_nodate}:

                \begin{lstlisting}[style=bash]
wget https://raw.githubusercontent.com/microsoft/Azure-Kinect-Sensor-SDK/\
develop/scripts/99-k4a.rules
sudo cp 99-k4a.rules /etc/udev/rules.d/
sudo udevadm control --reload-rules && sudo udevadm trigger
                \end{lstlisting}

            \item Außerhalb des Linux-Containers müssen die folgenden Einstellungen erfolgen, damit der unprivilegierte Linux-Container Zugriff auf die Azure Kinect erhält:
                \begin{lstlisting}[style=bash]
lxc config device add $CONTAINER_NAME microsoft_generic_superspeed_usb\
_hub usb vendorid=045e productid=097a
lxc config device add $CONTAINER_NAME microsoft_generic_usb\
_hub usb vendorid=045e productid=097b
lxc config device add $CONTAINER_NAME azure_kinect_depth\
_camera usb vendorid=045e productid=097c
lxc config device add $CONTAINER_NAME azure_kinect_4k\
_camera usb vendorid=045e productid=097d
lxc config device add $CONTAINER_NAME azure_kinect_microphone_\
array usb vendorid=045e productid=097e
                \end{lstlisting}

            \item Den Linux-Container neustarten.
        \end{enumerate}

    \item Azure Kinect SDKs installieren:
        \begin{enumerate}[label*=\arabic*.]
            \item Microsoft-Paketquelle hinzufügen \cite{microsoftazure-kinect-sensor-sdk_installation_nodate}:
                \begin{lstlisting}[style=bash]
sudo apt install curl
curl https://packages.microsoft.com/keys/microsoft.asc | sudo apt-key add -
sudo apt-add-repository https://packages.microsoft.com/ubuntu/18.04/prod
                \end{lstlisting}

            \item Installieren der Azure Kinect SDKs:
                \begin{lstlisting}[style=bash]
sudo apt install k4a-tools=1.3.*
sudo apt install libk4a1.3-dev
                \end{lstlisting}

                \begin{redbox}{Wichtig:}
                    Zur installierten Version vom \quoteMark{k4a-tools}-Paket muss die jeweilige Development-Version vom \quoteMark{libk4a}-Paket installiert werden.

                    \begin{lstlisting}[style=bash]
sudo apt install libk4a<major>.<minor>-dev
                    \end{lstlisting}

                    z.B.

                    \begin{lstlisting}[style=bash]
sudo apt install libk4a1.3-dev
                    \end{lstlisting}
                \end{redbox}

            \item Um Paketkollisionen beim Updaten zu verhindern sollte das Paket \quoteMark{k4a-tools} vom Updateprozess ausgeschlossen werden:
                \begin{lstlisting}[style=bash]
sudo apt-mark hold k4a-tools
sudo apt-mark showhold
                \end{lstlisting}

            \item Nun kann das Paket \quoteMark{libk4abt1.0-dev} ohne Paketkollisionen installiert werden:
                \begin{lstlisting}[style=bash]
sudo apt install libk4abt1.0-dev
                \end{lstlisting}
        \end{enumerate}

    \item Azure Kinect Firmware aktualisieren:
        \begin{enumerate}[label*=\arabic*.]
            \item Überprüfen ob die Azure Kinect erkannt wird:
                \begin{lstlisting}[style=bash]
AzureKinectFirmwareTool -l
                \end{lstlisting}

            \item Die Ausgabe sollte in etwa wie folgt aussehen:
                \begin{lstlisting}[style=bash]
 == Azure Kinect DK Firmware Tool ==
Found 1 connected devices:
1: Device "000846393812"
                \end{lstlisting}

            \item Die jeweilige Firmware von \href{https://github.com/microsoft/Azure-Kinect-Sensor-SDK/blob/develop/docs/usage.md#msis}{https://github.com/microsoft/Azure-Kinect-Sensor-SDK\\/blob/develop/docs/usage.md\#msis} herunterladen und flashen
                \begin{lstlisting}[style=bash]
wget https://download.microsoft.com/download/1/9/8/198048e8-63f2-45c6-\
8f96-1fd541d1b4bc/AzureKinectDK_Fw_1.6.102075014.bin
                \end{lstlisting}

            \item Die heruntergeladene Firmware flashen:
                \begin{lstlisting}[style=bash]
AzureKinectFirmwareTool -u AzureKinectDK_Fw_1.6.102075014.bin
                \end{lstlisting}
        \end{enumerate}

    \item Azure Kinect testen:
        \begin{enumerate}[label*=\arabic*.]
            \item Device-Stream testen:
                \begin{lstlisting}[style=bash]
k4aviewer
                \end{lstlisting}

            \item Body Tracking testen:\\
                Die OpenGL-Version, der Grafikkartenname und der eingesetzte Grafikkartentreiber kann mit den folgenden Befehlen abgefragt werden:
                \begin{lstlisting}[style=bash]
sudo apt install mesa-utils
glxinfo | grep -E "Device:|OpenGL version"
                \end{lstlisting}

                Die Ausgabe kann z.B. wie folgt aussehen:
                \begin{lstlisting}[style=bash]
    Device: AMD Radeon RX 5700 XT (NAVI10, DRM 3.35.0, 5.4.0-39-generic,\
 LLVM 10.0.0) (0x731f)
OpenGL version string: 4.6 (Compatibility Profile) Mesa 20.2.0-devel (git\
-c0c03f4 2020-06-27 bionic-oibaf-ppa)
                \end{lstlisting}

                \begin{redbox}{Wichtig:}
                    Für das Body Tracking wird ein Grafikkartentreiber mit mindestens \quoteMark{OpenGL 4.4}- und \quoteMark{OpenGL Compatibility Context}-Unterstützung benötigt. Bei den Open-Source-Mesa-Treibern tritt lauf Stand 29. März 2020 mit dem \quoteMark{k4a-tools}-Paket in der Version 1.3 oder niedriger der Fehler\newline\newline
                    \quoteMark{Shader Error: 0:2(38): error: `gl\_FragColor' undeclared\\0:2(38): error: value of type vec4 cannot be assigned to variable of type error}\newline\newline
                    auf, welcher das Starten der \quoteMark{k4abt\_simple\_3d\_viewer}-Anwendung verhindert.

                    Hierzu muss die \quoteMark{simple\_3d\_viewer}-Anwendung mit dem in dieser Arbeit erstellten GLSL-Patch\footnote{\href{https://github.com/microsoft/Azure-Kinect-Samples/commit/15c303fb9a3a6b3a2c54a57cb11996823389595a}{https://github.com/microsoft/Azure-Kinect-Samples/commit/15c303fb9a3a6b3a\\2c54a57cb11996823389595a}}, welcher bereits in den Hauptzweig des \quoteMark{Azure-Kinect-Samples}-Repositories eingepflegt wurde, neukompiliert werden. Getestet wurde das Kompilieren mit der Commit-Version \quoteMark{ac696e4015648dd82dd019e59d94ba169f3c81aa}. Wenn ein Fehler beim Kompilieren auftritt, kann nach dem Klonen des Repositories explizit mit dem folgenden Befehl auf den in dieser Arbeit verwendeten Stand des Repositories zugegriffen werden.
                    \begin{lstlisting}[style=bash]
git checkout ac696e4015648dd82dd019e59d94ba169f3c81aa
                    \end{lstlisting}

                    Zum kompilieren der \quoteMark{simple\_3d\_viewer}-Anwendung müssen die folgenden Schritte eingehalten werden.
                    \newcounter{bt_wichtig}
                    \begin{enumerate}[label=\arabic*.]
                        \item Das \quoteMark{Azure-Kinect-Samples}-Repository klonen:
                            \begin{lstlisting}[style=bash]
git clone --recursive https://github.com/microsoft/Azure\
-Kinect-Samples.git
cd Azure-Kinect-Samples/
                            \end{lstlisting}

                        \item Um einen Kompilierfehler\footnote{\label{fn:pull_request_54}\href{https://github.com/microsoft/Azure-Kinect-Samples/pull/54}{https://github.com/microsoft/Azure-Kinect-Samples/pull/54}} zu vermeiden muss in der Datei \quoteMark{CMakeLists.txt} zwischen \quoteMark{FIND\_PACKAGE(k4a REQUIRED)} und \quoteMark{FIND\_PACKAGE(k4abt REQUIRED)} der Text \quoteMark{FIND\_PACKAGE(k4arecord REQUIRED)} eingefügt werden. Um einen weiteren Kompilierfehler\footref{fn:pull_request_54} zu vermeiden muss zudem in der Datei \quoteMark{CMakeLists.txt} im Verzeichnis \quoteMark{body-tracking-samples/simple\_3d\_viewer/} zwischen \quoteMark{k4abt} und \quoteMark{window\_controller\_3d::window\_controller\_3d} der Text \quoteMark{k4arecord} eingefügt werden.
                        \setcounter{bt_wichtig}{\value{enumiii}}
                    \end{enumerate}
                \end{redbox}

                \begin{redbox}{}
                    \begin{enumerate}[label=\arabic*.]
                        \setcounter{enumiii}{\value{bt_wichtig}}
                        \item Notwendige Pakete zum Kompilieren der \quoteMark{simple\_3d\_viewer}-Anwendung installieren:
                            \begin{lstlisting}[style=bash]
sudo apt install libxi-dev cmake ninja-build
sudo apt install libxrandr-dev libxinerama-dev libxcursor\
-dev libgl-dev # for GLFW
sudo apt install mesa-vulkan-drivers # for Vulkan support in
                                         # GLFW (optional)
                            \end{lstlisting}

                        \item Die \quoteMark{simple\_3d\_viewer}-Anwendung komplieren und ausführbar machen:
                            \begin{lstlisting}[style=bash]
mkdir build
cd build
cmake .. -GNinja
ninja

chmod +x ./bin/simple_3d_viewer
                            \end{lstlisting}
                    \end{enumerate}
                \end{redbox}
%nm -gD /usr/lib/x86_64-linux-gnu/libk4arecord.so | grep -i "k4a_playback"

                Wenn eine Nvidia-GPU verbaut ist, kann die GPU-Beschleunigung eingesetzt werden. Hierzu muss das CUDA-Paket installiert werden:
                \begin{lstlisting}[style=bash]
k4abt_simple_3d_viewer
                \end{lstlisting}

                Daraufhin kann die \quoteMark{simple\_3d\_viewer}-Anwendung mit GPU-Beschleunigung gestartet werden:
                \begin{lstlisting}[style=bash]
k4abt_simple_3d_viewer
                \end{lstlisting}

                bzw.

                \begin{lstlisting}[style=bash]
~/Azure-Kinect-Samples/build/bin/simple_3d_viewer
                \end{lstlisting}

                Wenn keine Nvidia-GPU verbaut ist muss ansonsten der CPU-Modus verwendet werden:
                \begin{lstlisting}[style=bash]
k4abt_simple_3d_viewer CPU
                \end{lstlisting}

                bzw.

                \begin{lstlisting}[style=bash]
~/Azure-Kinect-Samples/build/bin/simple_3d_viewer CPU
                \end{lstlisting}
        \end{enumerate}

    \item Zum Builden der Roboter-Gesten-Anwendung müssen die folgenden Pakete installiert werden:
        \begin{lstlisting}[style=bash]
sudo apt install cmake gdb gcc python python-catkin-tools
        \end{lstlisting}

    \item Source Code kompilieren:\\
        \begin{redbox}{Wichtig:}
            Die Roboter-Gesten-Anwendung wird mit der Option \quoteMark{Release} kompiliert. Dies ist empfehlenswert, wenn die Roboter-Gesten-Anwendung im Produktivbetrieb eingesetzt werden soll. Die Option \quoteMark{Release} kann durch die Option \quoteMark{Debug} ersetzt werden um eine Entwicklungsversion mit zusätzlicher Ausgabe zu erstellen. Um eine Version zur Durchführung der in dieser Arbeit erläuterten Messungen zu erhalten wird die Option \quoteMark{Release} und zusätzlich die folgende Option verwendet:

            \begin{lstlisting}[style=bash]
--cmake-args -DCMAKE_CXX_FLAGS="-DMEASUREMENT"
            \end{lstlisting}
        \end{redbox}
        \begin{enumerate}[label*=\arabic*.]
            \item InterbotiX SDK kompilieren:
                \begin{lstlisting}[style=bash]
cd ~/interbotix_ws
rm -rf build devel/
catkin build -DCMAKE_BUILD_TYPE=Release
source ~/.bashrc
                \end{lstlisting}

            \item Die Roboter-Gesten-Anwendung aus dem Ordner \quoteMark{Repositories/Teach-Industrial-Robots} auf der CD/ISO klonen, Abhängigkeiten installieren und anschließend die Roboter-Gesten-Anwendung kompilieren:
                \begin{lstlisting}[style=bash]
mkdir -p ~/tir_ws/src
cd ~/tir_ws/
catkin build -DCMAKE_BUILD_TYPE=Release
echo "# ROS
source ~/tir_ws/devel/setup.bash" >> ~/.bashrc
source ~/.bashrc

cd ~/tir_ws/src
git clone <path to repository on cd or iso/Repositories/Teach-Industrial\
-Robots>

sudo apt install libglm-dev libeigen3-dev rapidjson-dev libk4a1.3\
-dev libk4abt1.0-dev

cd ~/tir_ws
catkin build -DCMAKE_BUILD_TYPE=Release
source ~/.bashrc
                \end{lstlisting}
        \end{enumerate}

    \item Entwicklungsumgebung installieren (optional):
        \begin{enumerate}[label*=\arabic*.]
            \item Temporärer Ordner erstellen:
                \begin{lstlisting}[style=bash]
echo "# Temporary user folder
sudo mkdir -p /run/user/$(id --user)
sudo chown $(whoami) /run/user/$(id --user)" >> ~/.bashrc
source ~/.bashrc
                \end{lstlisting}

            \item VS Code installieren:
                \begin{lstlisting}[style=bash]
sudo snap install code --classic
                \end{lstlisting}

            \item Starten der Entwicklungsumgebung mit dem folgenden Befehl:
                \begin{lstlisting}[style=bash]
code
                \end{lstlisting}

            \item Mittels Strg + P in VS Code das \quoteMark{Quick Open}-Menü öffnen und jeweils eine der folgenden Befehle eingeben um die notwendigen Plugins zu installieren:
                \begin{lstlisting}[style=bash]
ext install ms-vscode.cpptools
ext install ms-vscode.cmake-tools
ext install ms-python.python
ext install ms-iot.vscode-ros
                \end{lstlisting}
        \end{enumerate}

        \item Wenn der Speicherplatz des Linux-Containers zu knapp wird, kann der Speicherplatz des Linux-Containers wie folgt auf dem Host-System vergrößert werden:
            \begin{enumerate}[label*=\arabic*.]
                \item Image des Linux-Containers z.B. um 5 GB vergrößern:
                    \begin{lstlisting}[style=bash]
sudo truncate -s +5G /var/snap/lxd/common/lxd/disks/default.img
                    \end{lstlisting}

                \item Host-System neustarten

                \item Anschließend das Image des Linux-Containers einhängen und daraufhin die neue Größe des Images eingeben:
                    \begin{lstlisting}[style=bash]
sudo mount /var/snap/lxd/common/lxd/disks/default.img /<mount point>
sudo btrfs filesystem resize maxG /<mount point>
                    \end{lstlisting}
            \end{enumerate}
\end{enumerate}


\subsection{Starten der Roboter-Gesten-Anwendung}
Für die Roboter-Gesten-Anwendung wird als Scheduler SCHED\_DEADLINE eingesetzt, welcher für diese Arbeit feste Echtzeitfähigkeit im 50 ms Bereich garantieren kann.

\begin{figure}[htb]
	\centering
	\includegraphics[width=0.85\textwidth]{images/anhang/sched_deadline}
	\caption[Ablauf von SCHED\_DEADLINE]{Ablauf von SCHED\_DEADLINE\\Quelle: Eigene Darstellung, in Anlehnung an \cite{man_sched7_nodate}}
	\label{fig:sched_deadline}
\end{figure}
\FloatBarrier

Für den \quoteMark{SCHED\_DEADLINE}-Scheduler müssen hierfür die Parameter, welche in Abbildung \ref{fig:sched_deadline} ersichtlich sind, richtig eingestellt werden. Der \quoteMark{Runtime}-Parameter beträgt zur Einhaltung der 50 ms Frist hierbei 50000000 ns, der \quoteMark{Deadline}-Parameter beträgt 50000000 ns und der \quoteMark{Periode}-Parameter beträgt 50000000 ns. Der \quoteMark{Runtime}-Parameter steht hierbei für die maximal benötigte Ausführungszeit der Aufgabe, der \quoteMark{Deadline}-Parameter für wie lange genügend Zeit vorhanden ist um die Aufgabe auszuführen und der \quoteMark{Periode}-Parameter für wann die nächste Periode gestartet werden muss \cite{man_sched7_nodate}.\\

\begin{lstlisting}[style=bash, caption={Einstellungen des Start-Scripts der Roboter-Gesten-Anwendung}, label={lst:settings_for_start_script}]
USE_ROS_COMMUNICATION=true
USE_SIMULATION=true  # Only usable with "USE_ROS_COMMUNICATION=true"
LXC_USER_NAME="ubuntu"
LXC_INSTANCE="<container name>"
ADDITIONAL_APP_PARAMETERS="--move-home-at-exit --move-home-at-error"
TEACH_POSITIONS=true  # If "true" positions can be teached, Otherwise the positions
                       # in the file "robot_arm_positions.json" will be executed.
SCHED_RUNTIME=50000000
SCHED_DEADLINE=50000000
SCHED_PERIOD=50000000
\end{lstlisting}\leavevmode\newline

Die Roboter-Gesten-Anwendung kann mit dem Start-Script \quoteMark{start\_tir.sh} aus dem Ordner \quoteMark{Repositories/Teach-Industrial-Robots} auf der CD/ISO auf dem Host-System im Linux-Container gestartet und dem \quoteMark{SCHED\_DEADLINE}-Scheduler zugewiesen werden. Die einstellbaren Optionen des Start-Scripts sind im Listing \ref{lst:settings_for_start_script} ersichtlich. Zum Ausführen des Start-Scripts muss im Listing \ref{lst:settings_for_start_script} die Bash-Variable \quoteMark{LXC\_INSTANCE} mit dem Namen des Linux-Containers ersetzt werden. Die Bash-Variable \quoteMark{USE\_ROS\_COMMUNICATION} steuert ob die Roboter-Gesten-Anwendung über ROS kommunizieren oder eine direkt Verbindung mit dem Industrieroboter aufbauen soll. Um der Roboter-Gesten-Anwendung mitteilen zu können ob diese im Lern- oder Abspielmodus gestartet werden soll, wird die Bash-Variable \quoteMark{TEACH\_POSITIONS} eingesetzt. Mit der Bash-Variable \quoteMark{USE\_SIMULATION} kann eingestellt werden ob die Anwendung im Simulator gestartet werden soll. Um weitere von der Roboter-Gesten-Anwendung unterstützte Parameter anzuwenden kann die Bash-Variable \quoteMark{ADDITIONAL\_APP\_PARAMETERS} mit den folgenden Parametern modifiziert werden.

\begin{longtable}{|>{\raggedright\arraybackslash}p{0.2\textwidth}|>{\raggedright\arraybackslash}p{0.20\textwidth}|>{\raggedright\arraybackslash}p{0.50\textwidth}|}
\hline
\rowcolor{LightGray} \thead[c]{Kurzschreibweise\\des\\Parameters} & \thead[c]{Ausgeschriebener\\Parameter} & \thead[c]{Beschreibung}\\
\hline
-opf & \texttt{-{}-}overwrite-positions-file & Ermöglicht das Überschreiben der Positionsdatei. Zuvor sollte ein Backup der Positionsdatei gemacht werde um Datenverlust zu vermeiden.\\
\hline
-rrp & \texttt{-{}-}repeat-recorded-positions & Die aufgezeichneten Positionen aus der Positionsdatei wiederholen lassen\newline (Default-Positionsdatei: "\~{}/robot\_arm\_positions.json")\\
\hline
-tp & \texttt{-{}-}teach-positions & Positionen lernen\newline (Default-Positionsdatei: "\~{}/robot\_arm\_positions.json")\\
\hline
-rn="\phantom{}<name>" & \texttt{-{}-}robot-name="\phantom{}<name>" & Name des Roboters\newline (Default: "wx200")\\
\hline
-rm="\phantom{}<name>" & \texttt{-{}-}robot-model="\phantom{}<name>" & Robotermodell\newline (in der Regel der Name des Roboters\newline Default: "wx200")\\
\hline
-ur & \texttt{-{}-}use-ros & ROS oder eine direkte Verbindung zum Industrieroboter verwenden\\
\hline
-mhaex & \texttt{-{}-}move-home-at-exit & Beim Beenden der Anwendung in die Ausgangsposition zurückfahren\\
\hline
-mhaer & \texttt{-{}-}move-home-at-error & Bei einem Fehler der Anwendung in die Ausgangsposition zurückfahren\\
\hline
 & --help & Zeigt die Hilfe\\
\hline
\end{longtable}

%-----------------------------------------------

% Scheduler wechseln:
% cat /sys/block/sda/queue/scheduler

%-------------

% https://www.researchgate.net/publication/331290349_The_real-time_linux_kernel_A_survey_on_Preempt_RT

% -----------

% verwendeter Tiefensensor-Modus -> warum?

% Intel Real Sense Camera ZR300

%-------------------


% https://books.google.at/books?id=vVMoDwAAQBAJ&pg=PA229&lpg=PA229&dq=lxc+realtime&source=bl&ots=6aSRy33GKg&sig=ACfU3U13TeQVx0D1fZ5Z6DxLmKdeHTkXOg&hl=de&sa=X&ved=2ahUKEwj3mrWgyuvpAhUEGuwKHdcmDr8Q6AEwA3oECAoQAQ#v=onepage&q=lxc%20realtime&f=false
